\chapter{Lắng nghe trước}
\markboth{Lắng nghe trước}{Listen First}


\section{Thầy nghe điều em chưa nói thành lời.}
\begin{truyen}{Thầy nghe điều em chưa nói thành lời.}{Listen before advising.}
\chuhoa{C}{ó khi học sinh nản không vì lười, mà vì trong lòng có quá nhiều thứ chưa biết nói với ai. Một câu nghe trước, hỏi sau, đôi khi cứu cả một buổi học.}
\textit{(Sometimes students are discouraged not because they are lazy, but because they carry too much inside. One sentence that listens before advising can save an entire lesson.)}
\end{truyen}

\begin{baihoc}
Lắng nghe đúng lúc giúp học sinh mở lòng trước khi mở vở.
\textit{(Listening at the right time helps students open their hearts before they open their notebooks.)}
\end{baihoc}

\begin{ghinhoanh}
\textbf{Short English to remember:}

Listen first. Advice comes later.

\end{ghinhoanh}

\begin{tuhocnhanh}
1. Khi nản, em cần được nghe hay được dạy trước? \textit{(When discouraged, do you need to be heard or taught first?)}

2. Điều gì em đang giữ trong lòng mà ít ai biết? \textit{(What are you keeping inside that few people know?)}
\end{tuhocnhanh}
\ngancach


\section{Nếu em mệt, cứ nói thật với thầy.}
\begin{truyen}{Nếu em mệt, cứ nói thật với thầy.}{Honesty helps.}
\chuhoa{M}{ột học sinh có thể cố giấu sự mệt mỏi rất lâu. Nhưng khi được phép nói thật, các em bớt phải gồng và bắt đầu có cơ hội hồi lại.}
\textit{(A student can hide exhaustion for a very long time. But once allowed to tell the truth, they no longer need to pretend and can begin to recover.)}
\end{truyen}

\begin{baihoc}
Sự thật không làm học sinh yếu đi; nó giúp các em đỡ kiệt sức hơn.
\textit{(Truth does not weaken students; it helps them stop exhausting themselves.)}
\end{baihoc}

\begin{ghinhoanh}
\textbf{Short English to remember:}

Honesty is the first step back.

\end{ghinhoanh}

\begin{tuhocnhanh}
1. Em có dễ nói thật khi mình đang mệt không? \textit{(Is it easy for you to be honest when you are tired?)}

2. Điều gì khiến em ngại nói ra? \textit{(What makes you hesitate to say it out loud?)}
\end{tuhocnhanh}
\ngancach


\section{Thầy không nghe để phán xét, thầy nghe để hiểu.}
\begin{truyen}{Thầy không nghe để phán xét, thầy nghe để hiểu.}{Understand, not judge.}
\chuhoa{K}{hi học sinh cảm thấy mình sẽ bị quy lỗi ngay lập tức, các em chọn im lặng. Khi các em thấy mình được hiểu trước, các em mới dám thay đổi sau.}
\textit{(When students expect instant judgment, they choose silence. When they feel understood first, they dare to change later.)}
\end{truyen}

\begin{baihoc}
Phán xét làm học sinh khép lại; thấu hiểu cho các em một cánh cửa.
\textit{(Judgment makes students shut down; understanding gives them a door.)}
\end{baihoc}

\begin{ghinhoanh}
\textbf{Short English to remember:}

Understand first. Correct later.

\end{ghinhoanh}

\begin{tuhocnhanh}
1. Em phản ứng thế nào khi bị hiểu sai? \textit{(How do you react when you are misunderstood?)}

2. Em mong người lớn hiểu điều gì ở mình hơn? \textit{(What do you wish adults understood better about you?)}
\end{tuhocnhanh}
\ngancach


\section{Cứ nói từ đầu, đừng để mình mệt một mình.}
\begin{truyen}{Cứ nói từ đầu, đừng để mình mệt một mình.}{Don't carry it alone.}
\chuhoa{C}{ó những nỗi nản học bắt đầu rất nhỏ. Nhưng nếu để lâu không nói, nó lớn dần lên thành cảm giác muốn bỏ cả chặng đường.}
\textit{(Some discouragement begins very small. But when left unspoken, it grows into the feeling of wanting to quit the whole road.)}
\end{truyen}

\begin{baihoc}
Nói sớm là một cách tự cứu mình trước khi mọi thứ nặng hơn.
\textit{(Speaking up early is one way to save yourself before things grow heavier.)}
\end{baihoc}

\begin{ghinhoanh}
\textbf{Short English to remember:}

Speak early. Heal earlier.

\end{ghinhoanh}

\begin{tuhocnhanh}
1. Em thường im lặng đến mức nào rồi mới nói? \textit{(How long do you usually stay silent before speaking up?)}

2. Có chuyện gì em đáng ra nên nói sớm hơn? \textit{(Is there anything you should have said sooner?)}
\end{tuhocnhanh}
\ngancach


\section{Một câu hỏi thật đôi khi còn quý hơn một lời khuyên dài.}
\begin{truyen}{Một câu hỏi thật đôi khi còn quý hơn một lời khuyên dài.}{A true question helps.}
\chuhoa{K}{hi thầy hỏi đúng một câu, học sinh bắt đầu nghĩ lại và tự tìm lối ra. Có những lúc, câu hỏi đúng chữa nản tốt hơn bài giảng dài.}
\textit{(When a teacher asks the right question, students begin to rethink and find their own way out. Sometimes the right question heals discouragement better than a long lecture.)}
\end{truyen}

\begin{baihoc}
Câu hỏi đúng giúp học sinh tham gia vào việc tự cứu mình.
\textit{(The right question helps students take part in saving themselves.)}
\end{baihoc}

\begin{ghinhoanh}
\textbf{Short English to remember:}

Good questions wake up hope.

\end{ghinhoanh}

\begin{tuhocnhanh}
1. Câu hỏi nào khiến em suy nghĩ lâu nhất gần đây? \textit{(What question has made you think the longest lately?)}

2. Nếu tự hỏi mình một câu thật, em sẽ hỏi gì? \textit{(If you asked yourself one honest question, what would it be?)}
\end{tuhocnhanh}
\ngancach


\section{Thầy không cần em trả lời hay ngay; thầy cần em thật.}
\begin{truyen}{Thầy không cần em trả lời hay ngay; thầy cần em thật.}{Truth over polish.}
\chuhoa{N}{hiều học sinh nản vì nghĩ mình phải nói cho đúng, cho hay, cho giống người khác. Nhưng điều thầy cần trước tiên là sự thật để còn biết mà giúp.}
\textit{(Many students get discouraged because they think they must answer perfectly. But what a teacher first needs is the truth, so help becomes possible.)}
\end{truyen}

\begin{baihoc}
Sự chân thật là điểm bắt đầu của sự tiến bộ.
\textit{(Honesty is the starting point of progress.)}
\end{baihoc}

\begin{ghinhoanh}
\textbf{Short English to remember:}

Be real first. Improve later.

\end{ghinhoanh}

\begin{tuhocnhanh}
1. Em có hay cố tỏ ra ổn khi thật ra không ổn không? \textit{(Do you often pretend to be okay when you are not?)}

2. Điều gì làm em khó sống thật trong lớp? \textit{(What makes it hard for you to be honest in class?)}
\end{tuhocnhanh}
\ngancach


\section{Không sao nếu hôm nay em chỉ nói được một nửa.}
\begin{truyen}{Không sao nếu hôm nay em chỉ nói được một nửa.}{Half a step still counts.}
\chuhoa{K}{hông phải học sinh nào cũng đủ bình tĩnh để nói hết ngay lần đầu. Có khi chỉ cần các em nói được một nửa, thế là đủ để ngày mai tiến thêm một chút nữa.}
\textit{(Not every student is calm enough to say everything the first time. Sometimes speaking only half of it is already enough for tomorrow's next step.)}
\end{truyen}

\begin{baihoc}
Tiến bộ nhỏ vẫn là tiến bộ thật.
\textit{(Small progress is still real progress.)}
\end{baihoc}

\begin{ghinhoanh}
\textbf{Short English to remember:}

Half a step still matters.

\end{ghinhoanh}

\begin{tuhocnhanh}
1. Em có đang đòi mình phải hoàn hảo mới dám bắt đầu không? \textit{(Are you demanding perfection from yourself before starting?)}

2. Hôm nay em có thể nói thêm điều gì, dù chỉ một chút? \textit{(What can you say today, even if only a little more?)}
\end{tuhocnhanh}
\ngancach


\section{Thầy nhìn thấy sự cố gắng nhỏ mà em tưởng chẳng ai thấy.}
\begin{truyen}{Thầy nhìn thấy sự cố gắng nhỏ mà em tưởng chẳng ai thấy.}{Small effort matters.}
\chuhoa{N}{hiều học sinh bỏ cuộc vì nghĩ nỗ lực của mình quá nhỏ, không ai nhận ra. Nhưng một lời chỉ ra đúng phần tiến bộ nhỏ đó có thể cứu lại lòng tin của các em.}
\textit{(Many students quit because they think their effort is too small for anyone to notice. But one sentence naming that small progress can restore their confidence.)}
\end{truyen}

\begin{baihoc}
Được nhìn thấy đúng lúc khiến học sinh có lý do để đi tiếp.
\textit{(Being truly seen at the right time gives students a reason to keep going.)}
\end{baihoc}

\begin{ghinhoanh}
\textbf{Short English to remember:}

Notice small effort. It grows.

\end{ghinhoanh}

\begin{tuhocnhanh}
1. Lần gần nhất ai đó ghi nhận nỗ lực của em là khi nào? \textit{(When was the last time someone noticed your effort?)}

2. Em muốn người khác nhìn thấy điều gì ở mình hơn? \textit{(What do you wish others would notice more about you?)}
\end{tuhocnhanh}
\ngancach


\section{Học sinh cần một chỗ để thở trước khi cần một chỗ để thắng.}
\begin{truyen}{Học sinh cần một chỗ để thở trước khi cần một chỗ để thắng.}{Breathe before winning.}
\chuhoa{K}{hi áp lực đã đầy, học sinh không còn đủ chỗ trong lòng để học. Thầy giỏi đôi khi chỉ làm một việc rất nhỏ: cho các em một khoảng thở.}
\textit{(When pressure is already full, students no longer have room inside to learn. A good teacher sometimes does one small thing: giving them space to breathe.)}
\end{truyen}

\begin{baihoc}
Muốn học sinh mạnh lên, trước hết đừng để các em nghẹt thở.
\textit{(If you want students to become stronger, first do not let them suffocate.)}
\end{baihoc}

\begin{ghinhoanh}
\textbf{Short English to remember:}

Breathe first. Then begin again.

\end{ghinhoanh}

\begin{tuhocnhanh}
1. Điều gì làm em thấy nghẹt thở nhất khi học? \textit{(What makes you feel the most suffocated in study?)}

2. Em cần một khoảng thở kiểu nào để học tiếp? \textit{(What kind of breathing space do you need to keep learning?)}
\end{tuhocnhanh}
\ngancach


\section{Có những ngày thầy không kéo em đi nhanh hơn — thầy chỉ đi chậm lại cùng em.}
\begin{truyen}{Có những ngày thầy không kéo em đi nhanh hơn — thầy chỉ đi...}{Walk at their pace.}
\chuhoa{K}{hông phải lúc nào giúp học sinh cũng là thúc các em chạy nhanh. Có ngày, điều tử tế nhất là đi chậm lại cùng các em để các em không thấy mình bị bỏ phía sau.}
\textit{(Helping students is not always about pushing them faster. Some days, the kindest thing is simply slowing down beside them so they do not feel left behind.)}
\end{truyen}

\begin{baihoc}
Đi cùng học sinh đôi khi quan trọng hơn kéo học sinh.
\textit{(Walking with students is sometimes more important than pulling them.)}
\end{baihoc}

\begin{ghinhoanh}
\textbf{Short English to remember:}

Sometimes care means slowing down.

\end{ghinhoanh}

\begin{tuhocnhanh}
1. Em có từng thấy mình bị bỏ lại trong lớp không? \textit{(Have you ever felt left behind in class?)}

2. Nếu ai đó đi chậm lại cùng em, điều đó có ý nghĩa gì? \textit{(What would it mean if someone slowed down for you?)}
\end{tuhocnhanh}
\ngancach